% --- Archon physics formalization source begin ---
% archon:physics
% archon:covers PhyXMiniProblems/problem_phyx_mini_0394.lean
% archon:source-report reports/phyx_mini/problem_phyx_mini_0394.source.json

\chapter{Physics problem phyx\_mini\_0394}
\label{ch:PhyXMiniProblems_problem_phyx_mini_0394}

\paragraph{Problem source.}
\#\# Physical scenario

The main waterline into a tall building has a pressure of 600 kPa at 5-m elevation below ground level. The building is shown in the figure.

\#\# Question

How much extra pressure does a pump need to add to ensure a waterline pressure of 200 kPa at the top floor 150 m aboveground?

\#\# Answer choices

- A: A: 490kPa

- B: B: 1116kPa

- C: C: 154kPa

- D: D: 10.2kPa

\#\# Figure caption (auxiliary; use the image as primary evidence)

The image depicts a diagram of a building with the following elements:

- A tall building with multiple floors.
- Text labeling the "Top floor" and "Ground."
- A vertical pipe running from below ground level to the top floor of the building.
- At the bottom of the image, an underground "Pump" is shown connected to a "Water main."
- The diagram indicates that the building is 150 meters tall, and the water main is positioned 5 meters below ground.
- There is water flowing from the water main into the pump, and then upwards through the building.
- The letter "H" is marked at the top of the building on the pipe.

\#\# Dataset metadata

Laws of Thermodynamics; Physical Model Grounding Reasoning, Multi-Formula Reasoning

\paragraph{Recorded answer/context.}
B: B: 1116kPa

\paragraph{Figure/image path.}
/root/proposal\_for\_physic/science-mango/phyx\_mini\_run/phyx\_data/test\_image/394.png

\paragraph{Formalization target.}
create a compiling Lean file with sorry bodies at `PhyXMiniProblems/problem\_phyx\_mini\_0394.lean`. The Lean declarations must preserve the physical quantities, dimensions or dimensional roles, figure labels, governing-law hypotheses, and final relation expressed by this problem.
Use Mathlib/Physlib names found through LeanExplore where available. If a physics API is missing, introduce faithful local abstractions rather than scalar placeholder aliases.

\begin{theorem}[Physics formalization target]
\label{thm:physics:phyx_mini_0394:target}
\lean{PhyXMiniProblems.ProblemPhyXMini0394.problem_phyx_mini_0394}
\uses{def:physics:phyx-mini-0394:phyxminiproblems-problemphyxmini0394-lengthinmeters, def:physics:phyx-mini-0394:phyxminiproblems-problemphyxmini0394-pressureinpascals, def:physics:phyx-mini-0394:phyxminiproblems-problemphyxmini0394-tallbuildingwaterlinesetup, def:physics:phyx-mini-0394:phyxminiproblems-problemphyxmini0394-matchesstatedproblemdata, def:physics:phyx-mini-0394:phyxminiproblems-problemphyxmini0394-matchesprimarybuildingfigure, def:physics:phyx-mini-0394:phyxminiproblems-problemphyxmini0394-usesstandardwaterandgravity, def:physics:phyx-mini-0394:phyxminiproblems-problemphyxmini0394-hasphysicalwaterlineparameters, def:physics:phyx-mini-0394:phyxminiproblems-problemphyxmini0394-satisfiesfigureliftgeometry, def:physics:phyx-mini-0394:phyxminiproblems-problemphyxmini0394-satisfiesidealpumppressureriselaw, def:physics:phyx-mini-0394:phyxminiproblems-problemphyxmini0394-satisfieshydrostaticriserlaw, def:physics:phyx-mini-0394:phyxminiproblems-problemphyxmini0394-displayedpressurekilopascals, def:physics:phyx-mini-0394:phyxminiproblems-problemphyxmini0394-recordedanswerchoice, def:physics:phyx-mini-0394:phyxminiproblems-problemphyxmini0394-roundstonearestkilopascal}
The assigned autoformalize agent should translate this physics problem into Lean declarations in the covered file, with theorem and lemma proof bodies written as `by sorry`.
\end{theorem}
\begin{proof}
This is an autoformalization task, not a proof task. Produce statements that can later be proved by the physics prover without weakening the physical model.
\end{proof}

% --- Archon declaration topology begin ---
\section{Declaration topology}
\begin{definition}[Length Quantity]
  \label{def:physics:phyx-mini-0394:phyxminiproblems-problemphyxmini0394-lengthquantity}
  \lean{PhyXMiniProblems.ProblemPhyXMini0394.LengthQuantity}
  A nonnegative physical length
\end{definition}
\begin{proof}
  This is the declared model definition of Length Quantity.
\end{proof}
\begin{definition}[Elevation Quantity]
  \label{def:physics:phyx-mini-0394:phyxminiproblems-problemphyxmini0394-elevationquantity}
  \lean{PhyXMiniProblems.ProblemPhyXMini0394.ElevationQuantity}
  A signed physical elevation relative to the ground datum
\end{definition}
\begin{proof}
  This is the declared model definition of Elevation Quantity.
\end{proof}
\begin{definition}[Mass Density Quantity]
  \label{def:physics:phyx-mini-0394:phyxminiproblems-problemphyxmini0394-massdensityquantity}
  \lean{PhyXMiniProblems.ProblemPhyXMini0394.MassDensityQuantity}
  A nonnegative mass density, with dimension mass per volume
\end{definition}
\begin{proof}
  This is the declared model definition of Mass Density Quantity.
\end{proof}
\begin{definition}[Acceleration Quantity]
  \label{def:physics:phyx-mini-0394:phyxminiproblems-problemphyxmini0394-accelerationquantity}
  \lean{PhyXMiniProblems.ProblemPhyXMini0394.AccelerationQuantity}
  A nonnegative acceleration magnitude
\end{definition}
\begin{proof}
  This is the declared model definition of Acceleration Quantity.
\end{proof}
\begin{definition}[Pressure Quantity]
  \label{def:physics:phyx-mini-0394:phyxminiproblems-problemphyxmini0394-pressurequantity}
  \lean{PhyXMiniProblems.ProblemPhyXMini0394.PressureQuantity}
  A physical pressure, using Physlib's pressure dimension
\end{definition}
\begin{proof}
  This is the declared model definition of Pressure Quantity.
\end{proof}
\begin{definition}[length In Meters]
  \label{def:physics:phyx-mini-0394:phyxminiproblems-problemphyxmini0394-lengthinmeters}
  \lean{PhyXMiniProblems.ProblemPhyXMini0394.lengthInMeters}
  \uses{def:physics:phyx-mini-0394:phyxminiproblems-problemphyxmini0394-lengthquantity}
  Coherent-SI metre readout of a nonnegative physical length
\end{definition}
\begin{proof}
  This is the declared model definition of length In Meters.
\end{proof}
\begin{definition}[elevation In Meters]
  \label{def:physics:phyx-mini-0394:phyxminiproblems-problemphyxmini0394-elevationinmeters}
  \lean{PhyXMiniProblems.ProblemPhyXMini0394.elevationInMeters}
  \uses{def:physics:phyx-mini-0394:phyxminiproblems-problemphyxmini0394-elevationquantity}
  Coherent-SI metre readout of a signed elevation
\end{definition}
\begin{proof}
  This is the declared model definition of elevation In Meters.
\end{proof}
\begin{definition}[density In Kilograms Per Cubic Meter]
  \label{def:physics:phyx-mini-0394:phyxminiproblems-problemphyxmini0394-densityinkilogramspercubicmeter}
  \lean{PhyXMiniProblems.ProblemPhyXMini0394.densityInKilogramsPerCubicMeter}
  \uses{def:physics:phyx-mini-0394:phyxminiproblems-problemphyxmini0394-massdensityquantity}
  Coherent-SI kilogram-per-cubic-metre readout of a mass density
\end{definition}
\begin{proof}
  This is the declared model definition of density In Kilograms Per Cubic Meter.
\end{proof}
\begin{definition}[acceleration In Meters Per Second Squared]
  \label{def:physics:phyx-mini-0394:phyxminiproblems-problemphyxmini0394-accelerationinmeterspersecondsquared}
  \lean{PhyXMiniProblems.ProblemPhyXMini0394.accelerationInMetersPerSecondSquared}
  \uses{def:physics:phyx-mini-0394:phyxminiproblems-problemphyxmini0394-accelerationquantity}
  Coherent-SI metre-per-second-squared readout of an acceleration
\end{definition}
\begin{proof}
  This is the declared model definition of acceleration In Meters Per Second Squared.
\end{proof}
\begin{definition}[pressure In Pascals]
  \label{def:physics:phyx-mini-0394:phyxminiproblems-problemphyxmini0394-pressureinpascals}
  \lean{PhyXMiniProblems.ProblemPhyXMini0394.pressureInPascals}
  \uses{def:physics:phyx-mini-0394:phyxminiproblems-problemphyxmini0394-pressurequantity}
  Coherent-SI pascal readout of a physical pressure
\end{definition}
\begin{proof}
  This is the declared model definition of pressure In Pascals.
\end{proof}
\begin{definition}[pressure In Kilopascals]
  \label{def:physics:phyx-mini-0394:phyxminiproblems-problemphyxmini0394-pressureinkilopascals}
  \lean{PhyXMiniProblems.ProblemPhyXMini0394.pressureInKilopascals}
  \uses{def:physics:phyx-mini-0394:phyxminiproblems-problemphyxmini0394-pressurequantity, def:physics:phyx-mini-0394:phyxminiproblems-problemphyxmini0394-pressureinpascals}
  Kilopascal readout used by the prose and displayed answer choices
\end{definition}
\begin{proof}
  This is the declared model definition of pressure In Kilopascals.
\end{proof}
\begin{definition}[Pipeline Fluid]
  \label{def:physics:phyx-mini-0394:phyxminiproblems-problemphyxmini0394-pipelinefluid}
  \lean{PhyXMiniProblems.ProblemPhyXMini0394.PipelineFluid}
  Fluid transported from the main to the top-floor outlet
\end{definition}
\begin{proof}
  This is the declared model definition of Pipeline Fluid.
\end{proof}
\begin{definition}[Pressure Reference]
  \label{def:physics:phyx-mini-0394:phyxminiproblems-problemphyxmini0394-pressurereference}
  \lean{PhyXMiniProblems.ProblemPhyXMini0394.PressureReference}
  Common reference convention for every line-pressure readout
\end{definition}
\begin{proof}
  This is the declared model definition of Pressure Reference.
\end{proof}
\begin{definition}[Hydraulic Location]
  \label{def:physics:phyx-mini-0394:phyxminiproblems-problemphyxmini0394-hydrauliclocation}
  \lean{PhyXMiniProblems.ProblemPhyXMini0394.HydraulicLocation}
  Hydraulic locations distinguished by the supplied building diagram
\end{definition}
\begin{proof}
  This is the declared model definition of Hydraulic Location.
\end{proof}
\begin{definition}[Figure Label]
  \label{def:physics:phyx-mini-0394:phyxminiproblems-problemphyxmini0394-figurelabel}
  \lean{PhyXMiniProblems.ProblemPhyXMini0394.FigureLabel}
  Text and symbol labels printed in the primary figure
\end{definition}
\begin{proof}
  This is the declared model definition of Figure Label.
\end{proof}
\begin{definition}[Pipe Segment]
  \label{def:physics:phyx-mini-0394:phyxminiproblems-problemphyxmini0394-pipesegment}
  \lean{PhyXMiniProblems.ProblemPhyXMini0394.PipeSegment}
  The two visible pipe connections in the primary figure
\end{definition}
\begin{proof}
  This is the declared model definition of Pipe Segment.
\end{proof}
\begin{definition}[Tall Building Waterline Figure]
  \label{def:physics:phyx-mini-0394:phyxminiproblems-problemphyxmini0394-tallbuildingwaterlinefigure}
  \lean{PhyXMiniProblems.ProblemPhyXMini0394.TallBuildingWaterlineFigure}
  \uses{def:physics:phyx-mini-0394:phyxminiproblems-problemphyxmini0394-lengthquantity, def:physics:phyx-mini-0394:phyxminiproblems-problemphyxmini0394-elevationquantity, def:physics:phyx-mini-0394:phyxminiproblems-problemphyxmini0394-hydrauliclocation, def:physics:phyx-mini-0394:phyxminiproblems-problemphyxmini0394-figurelabel, def:physics:phyx-mini-0394:phyxminiproblems-problemphyxmini0394-pipesegment}
  Structured transcription of image 394.png. Elevations and dimensioned markers are physical quantities; the Boolean fields retain qualitative visual evidence that is not needed in the final arithmetic
\end{definition}
\begin{proof}
  This is the declared model definition of Tall Building Waterline Figure.
\end{proof}
\begin{definition}[Waterline Regime]
  \label{def:physics:phyx-mini-0394:phyxminiproblems-problemphyxmini0394-waterlineregime}
  \lean{PhyXMiniProblems.ProblemPhyXMini0394.WaterlineRegime}
  Idealization used for the school-level pump calculation
\end{definition}
\begin{proof}
  This is the declared model definition of Waterline Regime.
\end{proof}
\begin{definition}[Tall Building Waterline Setup]
  \label{def:physics:phyx-mini-0394:phyxminiproblems-problemphyxmini0394-tallbuildingwaterlinesetup}
  \lean{PhyXMiniProblems.ProblemPhyXMini0394.TallBuildingWaterlineSetup}
  \uses{def:physics:phyx-mini-0394:phyxminiproblems-problemphyxmini0394-lengthquantity, def:physics:phyx-mini-0394:phyxminiproblems-problemphyxmini0394-massdensityquantity, def:physics:phyx-mini-0394:phyxminiproblems-problemphyxmini0394-accelerationquantity, def:physics:phyx-mini-0394:phyxminiproblems-problemphyxmini0394-pressurequantity, def:physics:phyx-mini-0394:phyxminiproblems-problemphyxmini0394-pipelinefluid, def:physics:phyx-mini-0394:phyxminiproblems-problemphyxmini0394-pressurereference, def:physics:phyx-mini-0394:phyxminiproblems-problemphyxmini0394-tallbuildingwaterlinefigure, def:physics:phyx-mini-0394:phyxminiproblems-problemphyxmini0394-waterlineregime}
  The building waterline and its independent physical observables. In particular, pumpAddedPressure is an independent pressure increment. It is not defined to be any answer-choice value; the governing laws below relate it to the main, discharge, and top-floor pressures
\end{definition}
\begin{proof}
  This is the declared model definition of Tall Building Waterline Setup.
\end{proof}
\begin{definition}[Matches Stated Problem Data]
  \label{def:physics:phyx-mini-0394:phyxminiproblems-problemphyxmini0394-matchesstatedproblemdata}
  \lean{PhyXMiniProblems.ProblemPhyXMini0394.MatchesStatedProblemData}
  \uses{def:physics:phyx-mini-0394:phyxminiproblems-problemphyxmini0394-pressureinkilopascals, def:physics:phyx-mini-0394:phyxminiproblems-problemphyxmini0394-tallbuildingwaterlinesetup}
  Pressure values stated in the prose and the named working fluid
\end{definition}
\begin{proof}
  This is the declared model definition of Matches Stated Problem Data.
\end{proof}
\begin{definition}[Matches Primary Building Figure]
  \label{def:physics:phyx-mini-0394:phyxminiproblems-problemphyxmini0394-matchesprimarybuildingfigure}
  \lean{PhyXMiniProblems.ProblemPhyXMini0394.MatchesPrimaryBuildingFigure}
  \uses{def:physics:phyx-mini-0394:phyxminiproblems-problemphyxmini0394-lengthinmeters, def:physics:phyx-mini-0394:phyxminiproblems-problemphyxmini0394-elevationinmeters, def:physics:phyx-mini-0394:phyxminiproblems-problemphyxmini0394-tallbuildingwaterlinesetup}
  Primary-image evidence: ground is the zero datum, the water main and pump are 5 m below it, and the outlet marked H is 150 m above it. The pipe runs from the main through the pump and upward to H
\end{definition}
\begin{proof}
  This is the declared model definition of Matches Primary Building Figure.
\end{proof}
\begin{definition}[Uses Standard Water And Gravity]
  \label{def:physics:phyx-mini-0394:phyxminiproblems-problemphyxmini0394-usesstandardwaterandgravity}
  \lean{PhyXMiniProblems.ProblemPhyXMini0394.UsesStandardWaterAndGravity}
  \uses{def:physics:phyx-mini-0394:phyxminiproblems-problemphyxmini0394-densityinkilogramspercubicmeter, def:physics:phyx-mini-0394:phyxminiproblems-problemphyxmini0394-accelerationinmeterspersecondsquared, def:physics:phyx-mini-0394:phyxminiproblems-problemphyxmini0394-tallbuildingwaterlinesetup}
  The standard room-temperature water density and school-level terrestrial gravity calibration that reproduce the precision of the recorded answer
\end{definition}
\begin{proof}
  This is the declared model definition of Uses Standard Water And Gravity.
\end{proof}
\begin{definition}[Has Physical Waterline Parameters]
  \label{def:physics:phyx-mini-0394:phyxminiproblems-problemphyxmini0394-hasphysicalwaterlineparameters}
  \lean{PhyXMiniProblems.ProblemPhyXMini0394.HasPhysicalWaterlineParameters}
  \uses{def:physics:phyx-mini-0394:phyxminiproblems-problemphyxmini0394-lengthinmeters, def:physics:phyx-mini-0394:phyxminiproblems-problemphyxmini0394-densityinkilogramspercubicmeter, def:physics:phyx-mini-0394:phyxminiproblems-problemphyxmini0394-accelerationinmeterspersecondsquared, def:physics:phyx-mini-0394:phyxminiproblems-problemphyxmini0394-pressureinpascals, def:physics:phyx-mini-0394:phyxminiproblems-problemphyxmini0394-tallbuildingwaterlinesetup}
  Positivity conditions expressing that the pump adds, rather than removes, pressure
\end{definition}
\begin{proof}
  This is the declared model definition of Has Physical Waterline Parameters.
\end{proof}
\begin{definition}[Satisfies Figure Lift Geometry]
  \label{def:physics:phyx-mini-0394:phyxminiproblems-problemphyxmini0394-satisfiesfigureliftgeometry}
  \lean{PhyXMiniProblems.ProblemPhyXMini0394.SatisfiesFigureLiftGeometry}
  \uses{def:physics:phyx-mini-0394:phyxminiproblems-problemphyxmini0394-lengthinmeters, def:physics:phyx-mini-0394:phyxminiproblems-problemphyxmini0394-elevationinmeters, def:physics:phyx-mini-0394:phyxminiproblems-problemphyxmini0394-tallbuildingwaterlinesetup}
  The stored lift is the elevation difference from the main to outlet H
\end{definition}
\begin{proof}
  This is the declared model definition of Satisfies Figure Lift Geometry.
\end{proof}
\begin{definition}[Satisfies Ideal Pump Pressure Rise Law]
  \label{def:physics:phyx-mini-0394:phyxminiproblems-problemphyxmini0394-satisfiesidealpumppressureriselaw}
  \lean{PhyXMiniProblems.ProblemPhyXMini0394.SatisfiesIdealPumpPressureRiseLaw}
  \uses{def:physics:phyx-mini-0394:phyxminiproblems-problemphyxmini0394-pressureinpascals, def:physics:phyx-mini-0394:phyxminiproblems-problemphyxmini0394-tallbuildingwaterlinesetup}
  The pump discharge pressure is the supply pressure plus the pump increment
\end{definition}
\begin{proof}
  This is the declared model definition of Satisfies Ideal Pump Pressure Rise Law.
\end{proof}
\begin{definition}[Satisfies Hydrostatic Riser Law]
  \label{def:physics:phyx-mini-0394:phyxminiproblems-problemphyxmini0394-satisfieshydrostaticriserlaw}
  \lean{PhyXMiniProblems.ProblemPhyXMini0394.SatisfiesHydrostaticRiserLaw}
  \uses{def:physics:phyx-mini-0394:phyxminiproblems-problemphyxmini0394-lengthinmeters, def:physics:phyx-mini-0394:phyxminiproblems-problemphyxmini0394-densityinkilogramspercubicmeter, def:physics:phyx-mini-0394:phyxminiproblems-problemphyxmini0394-accelerationinmeterspersecondsquared, def:physics:phyx-mini-0394:phyxminiproblems-problemphyxmini0394-pressureinpascals, def:physics:phyx-mini-0394:phyxminiproblems-problemphyxmini0394-tallbuildingwaterlinesetup}
  Hydrostatic pressure balance along the lossless vertical riser: pressure at the top equals discharge pressure minus rho * g * Δz
\end{definition}
\begin{proof}
  This is the declared model definition of Satisfies Hydrostatic Riser Law.
\end{proof}
\begin{lemma}[total Vertical Lift In Meters]
  \label{lem:physics:phyx-mini-0394:phyxminiproblems-problemphyxmini0394-totalverticalliftinmeters}
  \lean{PhyXMiniProblems.ProblemPhyXMini0394.totalVerticalLiftInMeters}
  \uses{def:physics:phyx-mini-0394:phyxminiproblems-problemphyxmini0394-lengthinmeters, def:physics:phyx-mini-0394:phyxminiproblems-problemphyxmini0394-tallbuildingwaterlinesetup, def:physics:phyx-mini-0394:phyxminiproblems-problemphyxmini0394-matchesprimarybuildingfigure, def:physics:phyx-mini-0394:phyxminiproblems-problemphyxmini0394-satisfiesfigureliftgeometry}
  The figure geometry gives a total rise of 150 - (-5) = 155 m
\end{lemma}
\begin{proof}
  The conclusion follows from the governing relations and figure readouts named in the statement of total Vertical Lift In Meters.
\end{proof}
\begin{lemma}[required Pump Added Pressure In Pascals]
  \label{lem:physics:phyx-mini-0394:phyxminiproblems-problemphyxmini0394-requiredpumpaddedpressureinpascals}
  \lean{PhyXMiniProblems.ProblemPhyXMini0394.requiredPumpAddedPressureInPascals}
  \uses{def:physics:phyx-mini-0394:phyxminiproblems-problemphyxmini0394-pressureinpascals, def:physics:phyx-mini-0394:phyxminiproblems-problemphyxmini0394-tallbuildingwaterlinesetup, def:physics:phyx-mini-0394:phyxminiproblems-problemphyxmini0394-matchesstatedproblemdata, def:physics:phyx-mini-0394:phyxminiproblems-problemphyxmini0394-matchesprimarybuildingfigure, def:physics:phyx-mini-0394:phyxminiproblems-problemphyxmini0394-usesstandardwaterandgravity, def:physics:phyx-mini-0394:phyxminiproblems-problemphyxmini0394-satisfiesfigureliftgeometry, def:physics:phyx-mini-0394:phyxminiproblems-problemphyxmini0394-satisfiesidealpumppressureriselaw, def:physics:phyx-mini-0394:phyxminiproblems-problemphyxmini0394-satisfieshydrostaticriserlaw}
  Substitution in the two pressure-balance laws gives the unrounded pump increment 1,115,962 Pa = 1,115.962 kPa
\end{lemma}
\begin{proof}
  The conclusion follows from the governing relations and figure readouts named in the statement of required Pump Added Pressure In Pascals.
\end{proof}
\begin{definition}[Answer Choice]
  \label{def:physics:phyx-mini-0394:phyxminiproblems-problemphyxmini0394-answerchoice}
  \lean{PhyXMiniProblems.ProblemPhyXMini0394.AnswerChoice}
  Labels of the four answer choices printed with the problem
\end{definition}
\begin{proof}
  This is the declared model definition of Answer Choice.
\end{proof}
\begin{definition}[displayed Pressure Kilopascals]
  \label{def:physics:phyx-mini-0394:phyxminiproblems-problemphyxmini0394-displayedpressurekilopascals}
  \lean{PhyXMiniProblems.ProblemPhyXMini0394.displayedPressureKilopascals}
  \uses{def:physics:phyx-mini-0394:phyxminiproblems-problemphyxmini0394-answerchoice}
  Kilopascal value printed beside each displayed answer label
\end{definition}
\begin{proof}
  This is the declared model definition of displayed Pressure Kilopascals.
\end{proof}
\begin{definition}[recorded Answer Choice]
  \label{def:physics:phyx-mini-0394:phyxminiproblems-problemphyxmini0394-recordedanswerchoice}
  \lean{PhyXMiniProblems.ProblemPhyXMini0394.recordedAnswerChoice}
  \uses{def:physics:phyx-mini-0394:phyxminiproblems-problemphyxmini0394-answerchoice}
  Answer label recorded in the dataset metadata
\end{definition}
\begin{proof}
  This is the declared model definition of recorded Answer Choice.
\end{proof}
\begin{definition}[Rounds To Nearest Kilopascal]
  \label{def:physics:phyx-mini-0394:phyxminiproblems-problemphyxmini0394-roundstonearestkilopascal}
  \lean{PhyXMiniProblems.ProblemPhyXMini0394.RoundsToNearestKilopascal}
  \uses{def:physics:phyx-mini-0394:phyxminiproblems-problemphyxmini0394-pressurequantity, def:physics:phyx-mini-0394:phyxminiproblems-problemphyxmini0394-pressureinkilopascals}
  A pressure readout rounds to a displayed whole-kilopascal value
\end{definition}
\begin{proof}
  This is the declared model definition of Rounds To Nearest Kilopascal.
\end{proof}
% --- Archon declaration topology end ---
% --- Archon physics formalization source end ---
